% Source: https://artofproblemsolving.com/community/c3h2228101_tricky_geo_problem
% Stable identity: aops-topic-2228101-q01
\textbf{The triangle.} By the law of sines, $\sin P=\dfrac{p}{2\cdot1}=\dfrac{\sqrt3}{2}$ and $\sin Q=\dfrac {q}{2}=\dfrac{1}{2}$, so $P\in\{60^\circ,120^\circ\}$ and $Q\in\{30^\circ,150^\circ\}$. Since $P+Q<180^\circ$, the options are $(60^\circ,30^\circ)$, which gives $R=90^\circ$ (excluded), and $(120^\circ,30^\circ)$. Hence $P=120^\circ$, $Q=R=30^\circ$ and $r=PQ=q=1$.

\textbf{Coordinates.} Put $P=(0,0)$, $Q=(1,0)$, $R=\left(-\dfrac{1}{2},\dfrac{\sqrt3}{2}\right)$ (so $PR=1$, $\angle P=120^\circ$, $QR=\sqrt3$). Then $S=\left(\dfrac{1}{2},0\right)$. Line $QR$ has direction $\left(-\dfrac{3}{2},\dfrac{\sqrt3}{2}\right)$, i.e. slope $-\dfrac{1}{\sqrt3}$, so the altitude from $P$ is $y=\sqrt3x$ and meets $QR$ at $E=\left(\dfrac{1}{4},\dfrac{\sqrt3}{4}\right)$ (the midpoint of $QR$, as $PQ=PR$).

\begin{asy}
usepackage("amsmath");
size(6cm);
defaultpen(fontsize(10pt));
pair P=(0,0), Q=(1,0), R=(-0.5,sqrt(3)/2);
pair Sm=(P+Q)/2; pair Ef=(0.25,sqrt(3)/4); pair O=(1/6,sqrt(3)/6);
draw(P--Q--R--cycle,linewidth(0.9));
draw(R--Sm,blue); draw(P--Ef,red);
fill(Sm--O--Ef--cycle,lightgray); draw(Sm--O--Ef--cycle,gray);
real m=0.05; pair u=unit(R-Q), v=unit(P-Ef);
draw(Ef+m*u--Ef+m*u+m*v--Ef+m*v);
dot(Sm);dot(Ef);dot(O);
label("$P$",P,SW);label("$Q$",Q,SE);label("$R$",R,N);
label("$S$",Sm,S);label("$E$",Ef,NE);label("$O$",O,W);
label("$1$",(P+R)/2,W);label("$\sqrt3$",R+0.75*(Q-R),NE);draw(arc(P,0.07,0,120));label("$120^\circ$",P+0.15*dir(100),fontsize(7pt));
\end{asy}

Line $RS$: from $S=\left(\dfrac{1}{2},0\right)$ in direction $R-S=\left(-1,\dfrac{\sqrt3}{2}\right)$, points $\left(\dfrac{1}{2}-t,\dfrac{\sqrt3}{2}t\right)$. On $y=\sqrt3x$: $\dfrac{\sqrt3}{2}t=\sqrt3\left(\dfrac{1}{2}-t\right)$, so $t=\dfrac{1}{3}$ and $O=\left(\dfrac{1}{6},\dfrac{\sqrt3}{6}\right)$.

\textbf{(a)} $RS=\sqrt{1+\dfrac{3}{4}}=\dfrac{\sqrt7}{2}$.

\textbf{(b)} By the shoelace formula with $S\left(\dfrac{1}{2},0\right)$, $O\left(\dfrac{1}{6},\dfrac{\sqrt3}{6}\right)$, $E\left(\dfrac{1}{4},\dfrac{\sqrt3}{4}\right)$:
\[[SOE]=\dfrac{1}{2}\left|\left(\dfrac{1}{6}-\dfrac{1}{2}\right)\dfrac{\sqrt3}{4}-\left(\dfrac{1}{4}-\dfrac{1}{2}\right)\dfrac{\sqrt3}{6}\right|=\dfrac{1}{2}\left|-\dfrac{\sqrt3}{12}+\dfrac{\sqrt3}{24}\right|=\dfrac{\sqrt3}{48}.\]

\textbf{(c)} $OE=\sqrt{\left(\dfrac{1}{12}\right)^2+\left(\dfrac{\sqrt3}{12}\right)^2}=\dfrac{2}{12}=\dfrac{1}{6}$.

\textbf{(d)} The area is $\dfrac{1}{2}\cdot1\cdot1\cdot\sin120^\circ=\dfrac{\sqrt3}{4}$ and the semiperimeter is $\dfrac{2+\sqrt3}{2}$, so the inradius is
\[\dfrac{\sqrt3}{4}\cdot\dfrac{2}{2+\sqrt3}=\dfrac{\sqrt3}{2(2+\sqrt3)}=\dfrac{\sqrt3(2-\sqrt3)}{2}=\sqrt3-\dfrac{3}{2} .\]

\textbf{Sum.}
\[\dfrac{\sqrt7}{2}+\dfrac{\sqrt3}{48}+\dfrac{1}{6}+\sqrt3-\dfrac{3}{2}=\boxed{\dfrac{24\sqrt7+49\sqrt3-64}{48}}\approx1.758 .\]
